% Nouveau EGA V (cote 25): archive cote 25, site page 12 (PDF page 13).
% typescript. Draft compared: none.
% AI-transcribed by Claude Opus 5.5 (claude-opus-5-5, adaptive thinking, effort high), Message Batches API, batch ega25b (msgbatch_01A7n5YuXubUPNK1zsepi8RG), 2026-11-13; not yet checked by eye.
\scanpage{25/12}
\mspage{\hand{9}}

\uline{Corollaire} 2.11. Soit $y\in Y_{\hand{\eta}}$ , pour que $Y_{\hand{\eta}}$ \var{\inserted{soit régulier (resp.}}{soit régulier (resp.}{insertion}{ega25_25_012-v0001} satisfasse la propriété
$R_k$ (réf) en y\hand{,} \var{\deleted{(}}{}{deletion}{ega25_25_012-v0002}resp. soit normal en y) il f et s que X satisfasse
la même propriété en x. \var{\xout{\unsure{Remarque}}}{}{typewriter strike-out}{ega25_25_012-v0003} Soit Z l'ensemble des points de X en
lesquels X n'est pas \var{\deleted{\var{\xout{\illegible{1}}}{}{typewriter strike-out}{ega25_25_012-v0005} (}}{}{deletion}{ega25_25_012-v0004}\var{\inserted{régulier (resp. \var{\deleted{$R_k$}}{}{deletion}{ega25_25_012-v0007} $R_k$}}{régulier (resp. \var{\deleted{$R_k$}}{}{deletion}{ega25_25_012-v0007} $R_k$}{insertion}{ega25_25_012-v0006}resp. \var{\deleted{n'est pas}}{}{deletion}{ega25_25_012-v0008} normal); pour que $Y_{\hand{\eta}}$ soit
\var{\inserted{régulier (resp.}}{régulier (resp.}{insertion}{ega25_25_012-v0009} $(R_k)$\hand{,} \var{\deleted{(}}{}{deletion}{ega25_25_012-v0010}resp. normal) il f et s que f(Z) soit fini i.e. dim $\overline{f(Z)}=0$.

\var{\margin{|}}{\placedmargin{|}}{marginal note placed}{ega25_25_012-v0011} Même démonstration que 2.8. et 2.9. Il faut donner les diverses
références assurant que Z est fermé (car il faut savoir qu'il est
constructible pour appliquer 2.3.).

\var{\xout{\unsure{Remarque 2.12}}}{}{typewriter strike-out}{ega25_25_012-v0012} On notera que dans 2.10. et 2.11. nous ne disons rien
sur les propriétés ``géométriques'' correspondantes, les résultats
envisagés étant de nature \hand{``}absolue\hand{''}. Nous Examinons maintenant les
propriétés de nature géométrique. (On pourrait si on le désire en pro-%
\var{\margin{||}}{\placedmargin{||}}{marginal note placed}{ega25_25_012-v0013}fiter pour changer de n$^{\text{o}}$).

\var{\inserted[typed]{Supposons que $f:X\hand{\to} P$ soit une \var{\deleted{immersion}}{}{deletion}{ega25_25_012-v0015}\var{\inserted{\unsure{unramifié|non ramifié}}}{\unsure{unramifié|non ramifié}}{insertion}{ega25_25_012-v0016}.}}{Supposons que $f:X\hand{\to} P$ soit une \var{\deleted{immersion}}{}{deletion}{ega25_25_012-v0015}\var{\inserted{\unsure{unramifié|non ramifié}}}{\unsure{unramifié|non ramifié}}{insertion}{ega25_25_012-v0016}.}{insertion}{ega25_25_012-v0014}
\uline{Théorème} 2.12. Soient $y\in Y_{\hand{\eta}}$ , x son image dans X. Pour que X soit
lisse sur k en x, il f et s que $Y_{\hand{\eta}}$ soit lisse sur $k(\hand{\eta})$ en y.

On peut supposer k algébriquement clos. Si $Y_{\hand{\eta}}$ est lisse sur $k(\hand{\eta})$
en y, il est régulier, donc comme $Y_{\hand{\eta}}$ estb plat sur X, X est régulier
en x (réf), donc il est lisse sur k en x puisque k est alg clos donc
parfait (réf). Pour la réciproque, on peut (quitte à remplacer X par
un voisinage ouvert de x) supposer X lisse, et (grâce au critère jaco-%
bien de lissité) défini dans un ouvert U de P par p équations
comme $X=V(f_1,..,f_p)$ , où les différentielles $df_1,..,df_p$ sont
partout linéairement indépendantes. Introduisantdes coordonnées affines
$S_1,..,S_n$ dans $\overset{\hand{\vee}}{P}$ et des coordonnées affines $T_1,..,T_n$ dans un voisinage
de x (par choix d'un hyperplan $H_\infty$ ne contenant pas x), $Y\hand{\hookrightarrow} U_{k(\hand{\eta})}$ est
alors égal à $V(f_1,..,f_p,\hand{\sum_i}S_iT_i-1\ )$ , et il suffit de vérifier que
les différentielles (relativement à $k(\hand{\eta})$ ) de $f_1,..,f_p$ ,$\hand{(\sum}S_iT_i\hand{)}-1$

\mspage{\hand{12}}
